In the present paper, we are going to develop a similar theory, in the case where the weak IMCF~\eqref{eq:IMCF} is replaced by the level set flow of weak solutions $w_p\in \CS^{1,\beta}_{\loc}(M \smallsetminus \Omega)$ to the boundary value problem
\begin{equation}
\label{eq:pb-intro}
 \begin{cases}
 \Delta_p w_p = \abs{\D w_p}^p &\text{on $M \smallsetminus \Int \Omega$,}\\
 w_p=0 & \text{on $\partial \Omega$,}\\
 w_p\to +\infty & \text{as $\dist(x,\partial \Omega) \to +\infty$.}
 \end{cases}
\end{equation}
The link between the above problem and the weak IMCF~\eqref{eq:IMCF} relies on the fact that $w_p\to w_1$ as $p\to 1^+$ locally uniformly on $M$~\cite{kotschwar_localgradientestimatesharmonic_2009,mari_flowlaplaceapproximationnew_2022, moser_inversemeancurvatureflow_2007,moser_inversemeancurvatureflow_2008}, provided some natural global requirements are met by the manifold $(M, g)$. On the other hand,
the solutions to problem~\eqref{eq:pb-intro} are directly connected to the notion of $p$-capacitary potential of a compact body $\Omega$. In fact, setting $w_p = -(p-1) \log u_p$ implies that $u_p$ is $p$-harmonic, that is $\Delta_p u_p = 0$. These relationships have been instrumental in demonstrating a series of geometric inequalities by means of monotonicity formulas, holding along the level sets of solutions to equation~\eqref{eq:pb-intro}. As $p \to 1^+$, these inequalities become increasingly close to the desired result.
This machinery, firstly introduced in the case $p=2$ for harmonic functions in~\cite{agostiniani_sharpgeometricinequalitiesclosed_2020, agostiniani_monotonicityformulaspotentialtheory_2020}, has proven to be powerful enough to produce an enhanced version of the Minkowski inequality~\cite{agostiniani_minkowskiinequalitiesnonlinearpotential_2022, fogagnolo_geometricaspectscapacitarypotentials_2019}, later extended to Riemannian manifolds with nonegative Ricci curvature~\cite{benatti_minkowskiinequalitycompleteriemannian_2022} as well as to the anisotropic setting~\cite{xia_anisotropiccapacityanisotropicminkowski_2022}.
 In~\cite{agostiniani_greenfunctionproofpositive_2021} and subsequently in~\cite{agostiniani_riemannianpenroseinequalitynonlinear_2022}, the authors used this approach on $3$-manifolds with nonnegative scalar curvature to prove the Riemannian Penrose inequality for a single black hole, based on the monotonic behaviour of a suitable $p$-harmonic version of the Hawking mass (see~\eqref{eq:p-Hawkingmass} below).
\emph{Throughout the manuscript, Riemannian manifolds are assumed to be smooth, connected, metrically complete, noncompact, with one single end.}


The main object of interest in the present paper is the following nonlinear potential theoretic version of Huisken's isoperimetric mass~\eqref{eq:isoperimetric_mass_intro}, that we call the {\em $p$-isocapacitary mass}. It involves the classical notion of $p$-capacity of a compact set $K \subset M$, that, in dimension $3$ and for $1 < p <3$, is given by
\begin{equation*}%\label{eq:p-capacity-intro}
 \ncapa_p(K) = \inf \set{\frac{1}{4\pi}\bigg(\frac{p-1}{3-p}\bigg)^{p-1}\int_{M\smallsetminus K} \abs{\D v}^p \dif \mu\st v \in \CS_c^\infty(M),\, v \geq 1\text{ on } K}.
\end{equation*}
When the boundary of $K$ is regular enough, such infimum is realized by the $p$-capacitary potential of $K$, i.e., the unique $p$-harmonic function $u_p$, equal to $1$ on $\partial K$ and vanishing at infinity.
\begin{Definition}[$p$-isocapacitary mass]%\label{def:iso-p-capacitary_mass}
Let $(M,g)$ be a Riemannian $3$-manifold with compact boundary, and let $1<p<3$. Given a closed bounded subset $\Omega\subset M$ containing $\partial M$ with $\CS^{1}$-boundary the quasi-local $p$-isocapacitary mass of $\Omega$ is defined as
\begin{equation*}%\label{eq:quasi-local_isopcapacitary_mass}
 \ma_{\iso}^{\tp}(\Omega)=\frac{1}{2p \pi \ncapa_p(\partial \Omega)^{\frac{2}{3-p}}}\left( \abs{\Omega} -\frac{4 \pi}{3} \ncapa_p (\partial \Omega)^\frac{3}{3-p}\right).
\end{equation*}
The $p$-isocapacitary mass $\ma^{\tp}_{\iso}$ of $(M,g)$ is defined as \begin{equation}\label{eq:isopcapacitary_mass}
 \ma_{\iso}^{\tp} = \sup_{(\Omega_j)_{j \in \N}} \limsup_{j \to +\infty} \ma_{\iso}^{\tp}(\partial \Omega_j),
\end{equation}
where the supremum is taken among all exhaustions $\set{\Omega_j}_{j\in \N}$.
\end{Definition}
As for the isoperimetric mass, the $p$-isocapacitary mass might be any number in $[-\infty,+\infty]$. The special and particularly relevant case of the $2$-isocapacitary mass has been recently introduced and studied by Jauregui~\cite{jauregui_admmasscapacityvolumedeficit_2020}.

A first natural and fundamental question about these newly introduced quantities is whether they are nonnegative, on the class of $3$-manifolds with nonnegative scalar curvature, where a~solution to~\eqref{eq:pb-intro} exists for any bounded $\Omega$ with regular boundary. The latter mentioned property will be referred to as the \textit{strong $p$-nonparabolicity} of the manifold $(M,g)$, a terminology that interpolates between the notion of strong nonparabolicity introduced by Ni~\cite{ni_meanvaluetheoremsmanifolds_2007} and the notion of strong $1$-nonparabolicity, which was employed in~\cite{benatti_isoperimetricriemannianpenroseinequality_2022} and recalled above.

Our first main result is a nonlinear potential-theoretic version of the Riemannian Penrose inequality, that, although not sharp, implies the positive mass theorem for the $p$-isocapacitary mass, with its associated rigidity statement. We prove its validity under the following asymptotic integral gradient estimate:
\bgroup\leqnmode
\begin{equation*}
\begin{minipage}[t]{.9585\textwidth}%
Given any $\Omega \subset M$ closed bounded subset with smooth and connected boundary homologous to $\partial M$, the function $w_p\in \CS^1_{\loc}\big(M \smallsetminus \overline{\Omega}\big)$ that solves~\eqref{eq:pb-intro} satisfies
\begin{gather*}
\int_{\partial \Omega_t} \abs{ \D w_p}^2 \dif \sigma = o\bigl(\ee^{t/(p-1)}\bigr)
\end{gather*}
as $t \to +\infty$ where $\Omega_t= \set{ w_p \leq t}$.
\end{minipage}\tag{$\dag$}\label{eq:energy_hp}
\end{equation*}
\egroup

\begin{Theorem}[$p$-isocapacitary Riemannian Penrose inequality]
\label{thm:pmt-intro}
 Let $(M,g)$ be a strongly $p$-nonparabolic Riemannian $3$-manifold, $1<p<3$, with nonnegative scalar curvature and with smooth, compact, connected, minimal, possibly empty boundary. Assume also that $(M,g)$ satisfies~\eqref{eq:energy_hp} and $H_2( M, \partial M; \Z)=\set{0}$. Then
 \begin{equation}
\label{eq:p-penroseintro-crociata}
\ncapa_p(\partial M)^{\frac{1}{3-p}}\leq \frac{5-p}{2}\ma^{\tp}_{\iso}.
\end{equation}
In particular, $\ma^{\tp}_{\iso}\geq 0$ and it vanishes if and only if $(M,g)$ is isometric to the flat $3$-dimensional Euclidean space.
\end{Theorem}
